\uuid{2wGK}
\titre{Matrice/Inversibilité}
\auteur{hanani ; mzari}
\organisation{exo7}
\licence{CC BY-NC-SA 4.0}
\datecreate{2020-01-20}
\niveau{L1}
\module{Algèbre}
\chapitre{Matrice}
\sousChapitre{Inverse, méthode de Gauss}
\difficulte{3}
\theme{Inverse d'une matrice, Matrice, Inversibilité}
\exo7qcmid{lille-2.matrices.031}

\contenu{
\begin{question}
\qclassification{108.05}
\qtags{motcle=Matrice/Inversibilité}
On considère $ M_n(\Rr)$ l'ensemble des matrices carrées d'ordre $n$ à coefficients réels et $A$, $B$ et $C$ trois matrices non nulles deux à deux distinctes telles que $AB=AC$. Quelles sont les assertions vraies ?
\begin{answers}
\bad{$B=C$}.
\bad{$A=0$}.
\good{$A$ n'est pas inversible}.
\bad{Le rang de $A$ est $n$}.
\end{answers}
\begin{explanations} On a : $A(B-C)=0$, $A \neq 0$ et $B-C\neq 0$ (le produit de deux matrices peut être nul sans que les
deux matrices soient nulles).
\vskip0mm
Si $A$ est inversible, alors il existe une matrice $D$ telle que $DA=I$, où $I$ est la matrice identité. Donc  $(DA)(B-C) =B-C$. Or $(DA)(B-C)=D(A(B-C))$ et $A(B-C)=0$, donc $B-C=0$, ce qui est absurde. On déduit que $A$ n'est pas inversible et donc le rang de $A$ est $<n$.
\end{explanations}
\end{question}
}
